\chapter{Protocolli di Misurazione dei Parametri
DEQ$^2$}\label{protocolli-di-misurazione-dei-parametri-dequxb2}

\begin{quote}
\textbf{Scopo dell'appendice.} Trasformare le sette variabili
dell'apparato \((s_{iii}, z, \delta, \sigma\) esogene;
\(\Phi, A, \Omega, \varepsilon_{\text{dis}}\) in \(\mathbb{V}^{(4)}\))
in \emph{quantità misurabili in modo riproducibile da terzi
indipendenti}, su base trimestrale o annuale, a partire da fonti dati
pubbliche. L'appendice è la \emph{base tecnica} per il repository GitHub
\texttt{bilanciere-quantistico/predizioni-deq2} del Cap. 11.4 e per la
dashboard SPC-DEQ del Cap. 10.
\end{quote}



\section{Principi del Protocollo}\label{principi-del-protocollo}

\Hypoth{} \textbf{Cinque principi operativi.}

\begin{enumerate}
\def\labelenumi{\arabic{enumi}.}
\tightlist
\item
  \textbf{Replicabilità}: ogni misurazione deve essere ricalcolabile da
  terzi a partire da fonti pubbliche;
\item
  \textbf{Trasparenza dei proxy}: dichiarazione esplicita di proxy
  primario, secondari e fallback;
\item
  \textbf{Gestione del missing data}: regole esplicite di imputazione e
  di esclusione;
\item
  \textbf{Aggiornamento periodico}: ogni variabile ha frequenza di
  rilevazione fissata, con scadenze pubbliche;
\item
  \textbf{Intervallo di confidenza}: ogni stima è accompagnata da CI
  95\% via bootstrap.
\end{enumerate}

\Proved{} \textbf{Conseguenza.} I protocolli sotto consentono il calcolo
del punto
\(\boldsymbol{p}(t) = (\Phi, A, \Omega, \varepsilon_{\text{dis}})\) per
qualunque Stato membro ONU per cui esistano i dataset elencati. La
dashboard SPC-DEQ non richiede dati proprietari.



\section{\texorpdfstring{Protocollo per \(\Phi\) (Coerenza di
Fase)}{Protocollo per \textbackslash Phi (Coerenza di Fase)}}\label{protocollo-per-phi-coerenza-di-fase}

\Hypoth{} \textbf{Definizione operativa.} \(\Phi\) è il parametro
d'ordine di Kuramoto del campo \(\{\theta_j\}\) degli accoppiamenti
emittente-ricevente. Operativizzato come \emph{indice composito
multi-proxy}.

{\def\LTcaptype{none} % do not increment counter
{\small\begin{longtable}[]{@{}
  >{\raggedright\arraybackslash}p{(\linewidth - 8\tabcolsep) * \real{0.2000}}
  >{\raggedright\arraybackslash}p{(\linewidth - 8\tabcolsep) * \real{0.2000}}
  >{\raggedright\arraybackslash}p{(\linewidth - 8\tabcolsep) * \real{0.2000}}
  >{\raggedright\arraybackslash}p{(\linewidth - 8\tabcolsep) * \real{0.2000}}
  >{\raggedright\arraybackslash}p{(\linewidth - 8\tabcolsep) * \real{0.2000}}@{}}
\toprule\noalign{}
\begin{minipage}[b]{\linewidth}\raggedright
Componente
\end{minipage} & \begin{minipage}[b]{\linewidth}\raggedright
Proxy primario
\end{minipage} & \begin{minipage}[b]{\linewidth}\raggedright
Fonte
\end{minipage} & \begin{minipage}[b]{\linewidth}\raggedright
Frequenza
\end{minipage} & \begin{minipage}[b]{\linewidth}\raggedright
Peso
\end{minipage} \\
\midrule\noalign{}
\endhead
\bottomrule\noalign{}
\endlastfoot
\(\Phi^{\text{istituzionale}}\) & V-Dem Liberal Democracy Index &
\url{https://v-dem.net} & annuale (aprile) & 0.40 \\
\(\Phi^{\text{discorsivo}}\) & Pew Research \emph{common ground across
issues} (sei aree) & \url{https://pewresearch.org} & annuale & 0.25 \\
\(\Phi^{\text{fiduciario}}\) & Pew \emph{trust in federal government} &
\url{https://pewresearch.org} & annuale & 0.20 \\
\(\Phi^{\text{mediatico}}\) & Reuters Institute \emph{news trust index}
& \url{https://reutersinstitute.politics.ox.ac.uk} & annuale & 0.15 \\
\end{longtable}}
}

\Proved{} \textbf{Calcolo composito.} Ciascun sotto-indice è
normalizzato in \([0, 1]\) e combinato per media pesata:

\[
\Phi(t) = \sum_k w_k \cdot \tilde{\phi}_k(t)
\]

con \(\tilde{\phi}_k\) valore normalizzato e \(w_k\) peso da tabella.
Per Stati senza dataset Pew/Reuters, fallback a
\(\Phi^{\text{istituzionale}}\) con riduzione del peso totale
(rinormalizzazione).

\Hypoth{} \textbf{Pseudocodice (Python).}

\begin{verbatim}
def phi(country, year):
    vdem = fetch_vdem_lib_dem(country, year)
    pew_cg = fetch_pew_common_ground(country, year)
    pew_tr = fetch_pew_trust(country, year)
    rt_news = fetch_reuters_news_trust(country, year)
    weights = [0.40, 0.25, 0.20, 0.15]
    values = [vdem, pew_cg, pew_tr, rt_news]
    available = [v for v in values if v is not None]
    aw = [w for v, w in zip(values, weights) if v is not None]
    return sum(v*w for v, w in zip(available, aw)) / sum(aw)
\end{verbatim}

\Conj{} \textbf{Caveat.} Lo sbilanciamento Pew/Reuters verso paesi
anglofoni è una limitazione nota. Per Stati non-OECD, sostituire con
World Values Survey + Edelman Trust Barometer +
Afrobarometer/Latinobarómetro/Asiabarometer regionali.



\section{\texorpdfstring{Protocollo per \(A\)
(Antifragilità)}{Protocollo per A (Antifragilità)}}\label{protocollo-per-a-antifragilituxe0}

\Hypoth{} \textbf{Definizione operativa.}
\(A = \partial^2\Pi/\partial\sigma^2\) misurata via \emph{stress-test
retrospettivo} su tre crisi recenti: 2008-2009 (finanziaria), 2020
(pandemica), 2022-2023 (energetica/inflazionistica). Per ciascuna crisi:

\begin{itemize}
\tightlist
\item
  \(\Pi^{\text{baseline}}\): produttività totale media nei 3 anni
  pre-crisi;
\item
  \(\Pi^{\text{stress}}\): produttività media nei 3 anni post-crisi;
\item
  \(\sigma^{\text{shock}}\): deviazione standard di GDP, equity index,
  currency nei 12 mesi della crisi acuta.
\end{itemize}

\Hypoth{} \textbf{Stima.} Per ciascuna crisi \(i\):

\[
A_i = \frac{\Pi^{\text{stress}}_i - \Pi^{\text{baseline}}_i}{(\sigma^{\text{shock}}_i)^2}
\]

\(A^{\text{aggregato}} = \text{mediana}(\{A_1, A_2, A_3\})\) (più
robusto della media a outlier).

{\def\LTcaptype{none} % do not increment counter
{\small\begin{longtable}[]{@{}
  >{\raggedright\arraybackslash}p{(\linewidth - 2\tabcolsep) * \real{0.5000}}
  >{\raggedright\arraybackslash}p{(\linewidth - 2\tabcolsep) * \real{0.5000}}@{}}
\toprule\noalign{}
\begin{minipage}[b]{\linewidth}\raggedright
Fonte primaria
\end{minipage} & \begin{minipage}[b]{\linewidth}\raggedright
Per
\end{minipage} \\
\midrule\noalign{}
\endhead
\bottomrule\noalign{}
\endlastfoot
World Bank WDI (Total Factor Productivity) & \(\Pi\) \\
IMF International Financial Statistics & \(\sigma\) (volatilità
currency, equity) \\
OECD Productivity Database & TFP cross-country \\
\end{longtable}}
}

\Conj{} \textbf{Caveat.} \(A\) è la variabile più difficile in real-time
perché richiede \emph{eventi di stress} per manifestarsi. Per stime in
periodi senza crisi recenti, usare \emph{stress-test simulati} (Monte
Carlo controfattuali su shock di \(\sigma\) standardizzato) --- vedi
protocollo \S{}C.7.

\Hypoth{} \textbf{Pseudocodice.}

\begin{verbatim}
def A(country, crises=[(2008, 2009), (2020, 2020), (2022, 2023)]):
    A_values = []
    for c in crises:
        pre = mean_TFP(country, c[0]-3, c[0]-1)
        post = mean_TFP(country, c[1]+1, c[1]+3)
        shock = std_GDP(country, c[0], c[1])
        A_values.append((post - pre) / shock**2)
    return median(A_values)
\end{verbatim}



\section{\texorpdfstring{Protocollo per \(\Omega\)
(Involuzione)}{Protocollo per \textbackslash Omega (Involuzione)}}\label{protocollo-per-omega-involuzione}

\Hypoth{} \textbf{Definizione operativa.}
\(\Omega = \Delta E_{\text{dissipata}}/\Delta E_{\text{utile}} - 1\).
Operativizzato come \emph{complemento della Total Factor Productivity in
declino persistente}, con correzione per settori specifici dominanti.

{\def\LTcaptype{none} % do not increment counter
{\small\begin{longtable}[]{@{}
  >{\raggedright\arraybackslash}p{(\linewidth - 6\tabcolsep) * \real{0.2500}}
  >{\raggedright\arraybackslash}p{(\linewidth - 6\tabcolsep) * \real{0.2500}}
  >{\raggedright\arraybackslash}p{(\linewidth - 6\tabcolsep) * \real{0.2500}}
  >{\raggedright\arraybackslash}p{(\linewidth - 6\tabcolsep) * \real{0.2500}}@{}}
\toprule\noalign{}
\begin{minipage}[b]{\linewidth}\raggedright
Componente
\end{minipage} & \begin{minipage}[b]{\linewidth}\raggedright
Proxy
\end{minipage} & \begin{minipage}[b]{\linewidth}\raggedright
Fonte
\end{minipage} & \begin{minipage}[b]{\linewidth}\raggedright
Frequenza
\end{minipage} \\
\midrule\noalign{}
\endhead
\bottomrule\noalign{}
\endlastfoot
\(\Omega^{\text{TFP}}\) & \(-\Delta(\ln \text{TFP})/\Delta t\) in
finestra mobile 8 anni & OECD Productivity Database & annuale \\
\(\Omega^{\text{settore-critico}}\) & per Cina: contrazione property +
PPI; per USA: budget federale primario; per Brasile: deficit primario &
NBS (CN), BEA (US), IBGE (BR) & trimestrale \\
\(\Omega^{\text{governance}}\) & costo medio di rinegoziazione
legislativa (proxy: numero di vetoes, cloture votes, judicial overrides)
& sistemi legislativi nazionali & annuale \\
\end{longtable}}
}

\Proved{} \textbf{Combinazione.}
\(\Omega = \max(\Omega^{\text{TFP}}, \Omega^{\text{settore}}, \Omega^{\text{governance}})\)
--- l'involuzione \emph{si sente} dove è più alta, non dove è media.

\Hypoth{} \textbf{Esempio applicativo (Cina 2025).}
\(\Omega^{\text{TFP}} \approx 0{,}3\) (declino lento),
\(\Omega^{\text{settore}} \approx 1{,}2\) (property + deflazione PPI 10
trim.), \(\Omega^{\text{governance}} \approx 0{,}1\) (nessun vincolo
legislativo). \(\Omega^{\text{Cina}}_{2025} = 1{,}2\). Coerente con
stima Cap. 8 e Cap. 5.2.



\section{\texorpdfstring{Protocollo per \(\varepsilon_{\text{dis}}\)
(Disimpegno
Morale)}{Protocollo per \textbackslash varepsilon\_\{\textbackslash text\{dis\}\} (Disimpegno Morale)}}\label{protocollo-per-varepsilon_textdis-disimpegno-morale}

\Hypoth{} \textbf{Definizione operativa.}
\(\varepsilon_{\text{dis}} \in [0,1]\) è il coefficiente di disimpegno
morale collettivo (Bandura/Milgram). Indice composito.

{\def\LTcaptype{none} % do not increment counter
{\small\begin{longtable}[]{@{}
  >{\raggedright\arraybackslash}p{(\linewidth - 8\tabcolsep) * \real{0.2000}}
  >{\raggedright\arraybackslash}p{(\linewidth - 8\tabcolsep) * \real{0.2000}}
  >{\raggedright\arraybackslash}p{(\linewidth - 8\tabcolsep) * \real{0.2000}}
  >{\raggedright\arraybackslash}p{(\linewidth - 8\tabcolsep) * \real{0.2000}}
  >{\raggedright\arraybackslash}p{(\linewidth - 8\tabcolsep) * \real{0.2000}}@{}}
\toprule\noalign{}
\begin{minipage}[b]{\linewidth}\raggedright
Componente
\end{minipage} & \begin{minipage}[b]{\linewidth}\raggedright
Proxy
\end{minipage} & \begin{minipage}[b]{\linewidth}\raggedright
Fonte
\end{minipage} & \begin{minipage}[b]{\linewidth}\raggedright
Frequenza
\end{minipage} & \begin{minipage}[b]{\linewidth}\raggedright
Peso
\end{minipage} \\
\midrule\noalign{}
\endhead
\bottomrule\noalign{}
\endlastfoot
\(\varepsilon^{\text{partisan}}\) & Pew ``altro partito come minaccia
per la nazione'' & \url{https://pewresearch.org} & annuale & 0.30 \\
\(\varepsilon^{\text{turnout}}\) &
\(1 - \text{turnout elettorale normalizzato}\) & IDEA International &
per ciclo elettorale & 0.20 \\
\(\varepsilon^{\text{civic}}\) & \(1 - \text{volontariato pro capite}\)
& Charities Aid Foundation World Giving Index & annuale & 0.20 \\
\(\varepsilon^{\text{media-trust}}\) & \(1 - \text{trust mediatico}\) &
Edelman Trust Barometer & annuale & 0.15 \\
\(\varepsilon^{\text{intergroup}}\) & indice di dissociazione fra gruppi
(proxy: segregazione residenziale, omogamia educativa) & OECD Family
Database & quinquennale & 0.15 \\
\end{longtable}}
}

\Hypoth{} \textbf{Calcolo.} Identico a \(\Phi\) (media pesata,
normalizzazione, gestione missing data).



\section{\texorpdfstring{Protocolli per Variabili Esogene
\((s_{iii}, z, \delta, \sigma)\)}{Protocolli per Variabili Esogene (s\_\{iii\}, z, \textbackslash delta, \textbackslash sigma)}}\label{protocolli-per-variabili-esogene-s_iii-z-delta-sigma}

{\def\LTcaptype{none} % do not increment counter
{\small\begin{longtable}[]{@{}
  >{\raggedright\arraybackslash}p{(\linewidth - 6\tabcolsep) * \real{0.2500}}
  >{\raggedright\arraybackslash}p{(\linewidth - 6\tabcolsep) * \real{0.2500}}
  >{\raggedright\arraybackslash}p{(\linewidth - 6\tabcolsep) * \real{0.2500}}
  >{\raggedright\arraybackslash}p{(\linewidth - 6\tabcolsep) * \real{0.2500}}@{}}
\toprule\noalign{}
\begin{minipage}[b]{\linewidth}\raggedright
Variabile
\end{minipage} & \begin{minipage}[b]{\linewidth}\raggedright
Proxy primario
\end{minipage} & \begin{minipage}[b]{\linewidth}\raggedright
Fonte
\end{minipage} & \begin{minipage}[b]{\linewidth}\raggedright
Frequenza
\end{minipage} \\
\midrule\noalign{}
\endhead
\bottomrule\noalign{}
\endlastfoot
\(s_{iii}\) (legittimità etico-discorsiva) & Brand Finance Global Soft
Power Index &
\url{https://brandfinance.com/insights/global-soft-power-index} &
annuale \\
\(z\) (proiezione asse dialettico) & media ponderata Topografia \S{}5.1
(asse \(z\) del campo \(Q_0\)) & dataset corpus discorso pubblico &
annuale (proxy: indice complessità sintattica + diversità lessicale di
dichiarazioni esecutive) \\
\(\delta\) (orizzonte temporale) & OECD long-term planning indicator +
tasso di sconto sociale (Stern review) & OECD + Stern + AEI &
quinquennale \\
\(\sigma\) (volatilità ambientale) & std GDP + std CDS sovrani + std
petrolio + std eventi climatici estremi & IMF + Bloomberg + EM-DAT &
trimestrale \\
\end{longtable}}
}

\Conj{} \textbf{Caveat per \(z\).} La misurazione del \emph{contenuto
dialettico} del discorso pubblico richiede analisi NLP avanzata ---
pipeline non standardizzata. Proposta: usare il rapporto fra
\emph{propositions argumentate} e \emph{propositions assertive} in
corpus di dichiarazioni governative ufficiali (modello GPT-class come
strumento di codifica).



\section{\texorpdfstring{Stress-Test Simulato (per \(A\) in periodi
senza
crisi)}{Stress-Test Simulato (per A in periodi senza crisi)}}\label{stress-test-simulato-per-a-in-periodi-senza-crisi}

\Hypoth{} \textbf{Procedura Monte Carlo.} Quando la finestra delle tre
crisi recenti non è disponibile (paese giovane, dataset incompleto),
usare:

\begin{verbatim}
def A_simulated(country, year, n_simulations=5000):
    baseline = current_TFP(country, year)
    sigma_hist = historical_volatility(country, year)
    A_estimates = []
    for _ in range(n_simulations):
        shock = random_shock(distribution='lognormal', mean=0, sd=sigma_hist)
        # Modello dinamico: applicare shock, simulare risposta basata su:
        # - struttura federale (federalismo aumenta A)
        # - diversificazione produttiva (HHI inverso)
        # - opzionalità giuridica (numero di forme societarie)
        response = simulate_response(country, shock)
        A_estimates.append((response - baseline) / shock**2)
    return median(A_estimates), CI_95(A_estimates)
\end{verbatim}

\Proved{} \textbf{Validazione.} Calibrare il modello
\(simulate\_response\) su paesi con dataset crisi completo (USA, UE,
Giappone) e verificare che
\(A^{\text{simulato}} \approx A^{\text{empirico}}\) per quei paesi. Se
OK, applicare ad altri Stati.



\section{Tabella di Riepilogo
Operativo}\label{tabella-di-riepilogo-operativo}

{\def\LTcaptype{none} % do not increment counter
{\small\begin{longtable}[]{@{}
  >{\raggedright\arraybackslash}p{(\linewidth - 8\tabcolsep) * \real{0.2000}}
  >{\raggedright\arraybackslash}p{(\linewidth - 8\tabcolsep) * \real{0.2000}}
  >{\raggedright\arraybackslash}p{(\linewidth - 8\tabcolsep) * \real{0.2000}}
  >{\raggedright\arraybackslash}p{(\linewidth - 8\tabcolsep) * \real{0.2000}}
  >{\raggedright\arraybackslash}p{(\linewidth - 8\tabcolsep) * \real{0.2000}}@{}}
\toprule\noalign{}
\begin{minipage}[b]{\linewidth}\raggedright
Variabile
\end{minipage} & \begin{minipage}[b]{\linewidth}\raggedright
Frequenza calcolo
\end{minipage} & \begin{minipage}[b]{\linewidth}\raggedright
Tempo di latenza dati
\end{minipage} & \begin{minipage}[b]{\linewidth}\raggedright
CI tipica
\end{minipage} & \begin{minipage}[b]{\linewidth}\raggedright
Osservabilità del trend
\end{minipage} \\
\midrule\noalign{}
\endhead
\bottomrule\noalign{}
\endlastfoot
\(\Phi\) & annuale & 4-6 mesi (V-Dem) & \(\pm 0{,}03\) & media (segnali
a 1-2 anni) \\
\(A\) & post-stress + simulato annualmente & 3-5 anni (per stress reali)
& \(\pm 0{,}1\) & bassa (richiede crisi) \\
\(\Omega\) & trimestrale & 2-3 mesi & \(\pm 0{,}05\) & alta (segnali a
1-2 trimestri) \\
\(\varepsilon_{\text{dis}}\) & semestrale & 6-12 mesi & \(\pm 0{,}04\) &
media \\
\(s_{iii}\) & annuale & 3 mesi (Brand Finance) & \(\pm 1\) punto &
media \\
\(z\) & annuale & 6-12 mesi (NLP custom) & \(\pm 0{,}1\) & bassa \\
\(\delta\) & quinquennale & 1-3 anni & \(\pm 0{,}05\) & bassissima
(variabile slow) \\
\(\sigma\) & trimestrale & 1 mese & \(\pm 0{,}02\) & altissima \\
\end{longtable}}
}

\Hypoth{} \textbf{Conseguenza per la dashboard.} Il pannello aggregato
\(P_0 = d(\boldsymbol{p}, \mathcal{S})\) del Cap. 10 può essere
aggiornato \emph{trimestralmente} sulla base di
\(\Omega, \sigma, s_{iii}\) con interpolazione di
\(\Phi, \varepsilon_{\text{dis}}\) fra gli aggiornamenti annuali. La
risoluzione temporale effettiva del cruscotto è quindi \emph{quattro
volte/anno}.



\section{Limiti Espliciti del
Protocollo}\label{limiti-espliciti-del-protocollo}

\Conj{} \textbf{Cinque limiti riconosciuti.}

\begin{enumerate}
\def\labelenumi{\arabic{enumi}.}
\tightlist
\item
  \textbf{Bias verso paesi OECD}: la maggior parte dei proxy ha
  copertura completa solo per OECD + grandi emergenti. Per Stati piccoli
  o paesi a basso reddito: copertura ridotta, fallback a sotto-indici.
\item
  \textbf{Aggregazione vs sotto-attori}: la dashboard misura \emph{Stati
  nazionali}. Per sotto-attori (regioni, città, aziende) i protocolli
  vanno re-calibrati (vedi Cap. 13.5 vettore 2 di R\&D).
\item
  \textbf{Lag dati}: la dashboard rappresenta \emph{passato recente},
  non presente; per uso predittivo, applicare il sistema dinamico Cap.
  5.4 con \(\boldsymbol{p}(t-\Delta t)\) come stato iniziale.
\item
  \textbf{Sensibilità ai pesi}: i pesi nei sotto-indici sono scelte
  parametriche; il paper EN testerà la robustezza via sensitivity
  analysis.
\item
  \textbf{NLP per \(z\)}: la pipeline non è standardizzata; la
  riproducibilità di \(z\) è inferiore alle altre variabili.
\end{enumerate}

\Proved{} \textbf{Mitigazione.} Tutti e cinque i limiti sono
\emph{trasparenti} --- il repository GitHub riporta esplicitamente quali
paesi/anni hanno copertura completa vs parziale, con flag di
affidabilità.



\section{Sintesi dell'Appendice}\label{sintesi-dellappendice}

\Proved{} \textbf{Risultato.} Le sette variabili dell'apparato DEQ$^2$ sono
operativizzabili tramite indici compositi multi-proxy basati su dataset
pubblici (V-Dem, Pew, Reuters, OECD, IMF, Brand Finance, Edelman).
Frequenza di aggiornamento: trimestrale per \(\Omega\) e \(\sigma\),
semestrale per \(\varepsilon_{\text{dis}}\), annuale per
\(\Phi, s_{iii}, z, A\), quinquennale per \(\delta\). Latenza tipica
dato-stima: 1-12 mesi a seconda della variabile. Tutto il protocollo è
riproducibile in pseudocodice Python e implementabile come web service.

\Hypoth{} \textbf{Posizione strategica.} Senza un protocollo
riproducibile, le predizioni del Cap. 11 sarebbero tecnicamente non
falsificabili (le valutazioni dipenderebbero dalla scelta dei
misuratori). Con il protocollo C.1-C.9, le predizioni \emph{diventano}
falsificabili da terzi indipendenti, soddisfacendo il criterio (3) del
Cap. 11.0.



\emph{Appendice tecnica. Implementazione di riferimento: repository
\texttt{bilanciere-quantistico/predizioni-deq2} (in setup,
post-pubblicazione libro).}
